% --- INICIO DE PLANTILLA PERSONALIZADA ---
\documentclass[oneside,11pt]{Latex/Classes/PhDthesisPSnPDF}

% --- Paquetes Originales ---
\usepackage{blindtext}
\usepackage{actuarialangle} 
\usepackage{tabularx}
\usepackage{float}
\usepackage{multirow, array}
\usepackage[usenames,dvipsnames,svgnames,table]{xcolor}
\usepackage{longtable}
\usepackage{latexsym,amsmath,amssymb,amsfonts}
\usepackage{enumitem}
\usepackage{xparse}
\usepackage{booktabs}
\usepackage[round, sort, numbers]{natbib}  
\usepackage{adjustbox}
\usepackage[utf8]{inputenc}
\usepackage{caption}
\usepackage{graphicx}
\usepackage{hyperref}
\usepackage{cleveref}

% Definición de colores
\definecolor{miblue}{RGB}{68, 114, 196}

% --- CAMBIO CRÍTICO 1: Usar \input para macros en el preámbulo ---
% \include da problemas antes del begin document. \input es seguro.
\input{Latex/Macros/MacroFile1}

% --- Definiciones de seguridad (por si la clase falla al cargar alguna variable) ---
\providecommand{\facultad}[1]{\def\lafacultad{#1}}
\providecommand{\escudofacultad}[1]{\def\elescudo{#1}}
\providecommand{\degree}[1]{\def\elgrado{#1}}
\providecommand{\director}[1]{\def\eldirector{#1}}
\providecommand{\lugar}[1]{\def\ellugar{#1}}
\providecommand{\degreedate}[1]{\def\lafecha{#1}}

% --- Configuración para bloques de código de R (Pandoc) ---

%%%%%%%%%%%%%%%%%%%%%%%%%%%%%%%%%%%%%%%%%%%%%%%%%%%%%%%%%%%%%%%%%%%%%%%%%%%%%%%%
%                         DATOS (Conectados a RMarkdown)                       %
%%%%%%%%%%%%%%%%%%%%%%%%%%%%%%%%%%%%%%%%%%%%%%%%%%%%%%%%%%%%%%%%%%%%%%%%%%%%%%%%
\title{Análisis de factores demográficos y su impacto en los patrones de
consumo en una tienda minorista durante el año 2023.}
\author{Johan Infante \\ Eddie Palomino \\ Juan Covarrubia}

% Lógica para Facultad: Si la defines en el YAML la usa, si no, usa la default

\facultad{Facultad de Ciencias Económicas y Sociales \\ Escuela de
Estadistica y Ciencias Actuariales}
  \escudofacultad{Latex/Classes/Escudos/faces}

\degree{Computación I}
\director{Jesus Ochoa \\ Oliver Triveño}
\degreedate{Febrero 2026} 
\lugar{Caracas}

\portadatrue 

% Metadatos PDF
\keywords{}
\subject{}

% --- CORRECCIÓN PARA PANDOC NUEVO (Imágenes) ---
\makeatletter
\providecommand{\pandocbounded}[1]{#1}
\makeatother


%%%%%%%%%%%%%%%%%%%%%%%%%%%%%%%%%%%%%%%%%%%%%%%%%%%%%
%                   DOCUMENTO                       %
%%%%%%%%%%%%%%%%%%%%%%%%%%%%%%%%%%%%%%%%%%%%%%%%%%%%%
\begin{document}

\maketitle

%%%%%%%%%%%%%%%%%%%%%%%%%%%%%%%%%%%%%%%%%%%%%%%%%%%%%
%                  PRÓLOGO                          %
%%%%%%%%%%%%%%%%%%%%%%%%%%%%%%%%%%%%%%%%%%%%%%%%%%%%%
\frontmatter

% Puedes agregar dedicatorias desde el YAML si quieres, o dejarlas fijas aquí

%%%%%%%%%%%%%%%%%%%%%%%%%%%%%%%%%%%%%%%%%%%%%%%%%%%%%
%                   ÍNDICES                         %
%%%%%%%%%%%%%%%%%%%%%%%%%%%%%%%%%%%%%%%%%%%%%%%%%%%%%
\setcounter{secnumdepth}{3}
\setcounter{tocdepth}{3}

\tableofcontents



\cleardoublepage

  \cleardoublepage       % Empieza en página derecha (estándar en tesis)
  \chapter*{Resumen}     % Crea el título "Resumen" sin numeración (Capítulo *)
  \addcontentsline{toc}{chapter}{Resumen} % (Opcional) Lo añade al índice
  
  Resumen del trabajo de
investigación\ldots{}             % Aquí se pega el texto que escribiste en el YAML
  
  \cleardoublepage

%%%%%%%%%%%%%%%%%%%%%%%%%%%%%%%%%%%%%%%%%%%%%%%%%%%%%
%                   CONTENIDO (El Corazón)          %
%%%%%%%%%%%%%%%%%%%%%%%%%%%%%%%%%%%%%%%%%%%%%%%%%%%%%
\mainmatter
\def\baselinestretch{1.5}

% --- CAMBIO CRÍTICO 2: Aquí RMarkdown inyecta todo tu texto ---
\textbackslash chapter*\{Introducción\}

Introducción del trabajo de investigación

\chapter{Planteamiento del Problema}

En la actualidad, el sector minorista (retail) se enfrenta a un entorno
altamente competitivo donde la comprensión del cliente es vital. Tal
como lo demuestra el Estudio Retail 2024-2025 de la firma de análisis
GfK, las dinámicas de consumo y la frecuencia de compra están
directamente vinculadas a la segmentación generacional y el perfil del
comprador, obligando a las marcas a estudiar sus datos demográficos para
anticipar el comportamiento del mercado. Las empresas ya no solo venden
productos sino que gestionan información estratégica para optimizar su
inventario y fidelizar a sus consumidores.

A partir de esta realidad, surge la necesidad de aplicar este mismo
enfoque analítico a nivel micro. Al disponer de un conjunto de datos
transaccionales del año 2023 de una tienda \emph{retail}, existe la
oportunidad de extraer información valiosa que describa quiénes son los
clientes y cómo interactúan con los productos ofrecidos. El problema
central radica en transformar estos datos crudos en información
estructurada que permita comprender las dinámicas básicas de consumo sin
recurrir a modelos predictivos complejos, basándose estrictamente en el
comportamiento histórico registrado.

El problema central de esta investigación radica en determinar: ¿Existe
una asociación observable entre el perfil demográfico del consumidor y
sus hábitos de gasto y elección de productos en el conjunto de datos de
2023?

Argumento final. Debido a esto, se plantea las siguientes preguntas de
investigación:\\

\textbf{1.} ¿Cuál es el perfil demográfico de los clientes que conforman
el conjunto de datos de 2023 y cómo se distribuyen gráficamente según su
edad y género?

\textbf{2.} ¿Cómo fluctúa el comportamiento de las ventas y el gasto
promedio mensual a lo largo del año?

\textbf{3.} ¿Cuáles son las preferencias de compra en las distintas
categorías de productos según los diferentes grupos etarios de los
consumidores?\\

\section{Justificación}

La presente investigación se justifica desde un punto de vista práctico
y académico. En el ámbito práctico, el análisis descriptivo de bases de
datos de clientes permite a cualquier negocio identificar a su público
objetivo real y observar hechos concretos, como qué productos tienen
mayor salida en ciertos meses o qué rango de edad consume más tecnología
frente a artículos de hogar. Esta información es la piedra angular para
cualquier toma de decisiones comerciales, diseño de campañas o gestión
de inventario.

Desde el punto de vista académico, este trabajo representa una
aplicación directa y rigurosa de los fundamentos de estadística
descriptiva y computación del primer semestre. Permite demostrar el
dominio en la limpieza de datos, el uso de tablas de frecuencia, la
visualización de datos (como la construcción de pirámides poblacionales)
y el análisis de correlación simple. Al limitar el estudio a
herramientas descriptivas, se garantiza un manejo profundo y preciso de
los resultados obtenidos, sentando las bases analíticas necesarias para
futuras investigaciones más complejas.

\section{Objetivos}

\subsection{Objetivo General}
\begin{itemize}
  \item{nalizar el perfil demográfico y el comportamiento de compra de los clientes registrados en el conjunto de datos de retail durante el año 2023, utilizando técnicas de estadística descriptiva para la identificación de patrones de consumo.}
\end{itemize}
\subsection{Objetivos Específicos}
\begin{itemize}
  \item{}
  \item{}
  \item{}
  \item{}
\end{itemize}

\section{Cobertura}

\subsection{Horizontal}

Se limita al análisis de las variables demográficas (edad, género) y
transaccionales (categoría, cantidad, precio, total) contenidas en el
dataset de retail suministrado.

\begin{table}[ht]
\centering
\caption{Variables consideradas} 
\resizebox{10cm}{!} {
\begin{tabular}{ccc}
  \hline
\ \ \ \ \ NOMBRE DE LA VARIABLE \\ 
  \hline
   VARIABLE 1 \\ 
VARIABLE 2\\ 
VARIABLE n\\ 
   \hline
\end{tabular}
}
\end{table}

\subsection{Vertical}

1.3.2. Vertical: El estudio abarca desde el Análisis Exploratorio de
Datos (EDA) y estadística descriptiva, hasta la validación de hipótesis
mediante pruebas de independencia y la creación de un modelo de
regresión lineal simple.

\section{Periodo de Referencia}

Periodo de Referencia ``El análisis se basa en datos históricos
recolectados durante el año fiscal 2023, permitiendo observar el
comportamiento del consumidor a lo largo de los doce meses de dicho
periodo.''

\section{Variables}

Describir la variable en estudio en el trabajo de investigación

\chapter{Marco Teórico}

\section{Bases Teóricas}

Resumen de las bases teoricas explicadas en el capitulo y su
importancia\\

\subsection{Definición 1}

Descripción de la definición 1\\

\subsection{Definición 2}

Descripción de la definición 2\\

\subsection{Definición n}

Descripción de la definición n\\

\chapter{Marco Metódico}

En esta sección se presentan las metodologías y/o test necesarios
relacionados con la práctica divulgada en el marco teórico, así como las
fuentes de datos.

\section{Bases metódicas utilizadas en el trabajo de investigación}

Desarrollo\\

\subsection{Item 1}

Descripción\\

\subsection{Item 2}

Descripción\\

\subsubsection{Item 2.2}

Descripción\\

\subsection{Item n}

Descripción\\

\chapter{Análisis de Resultados}

\{\small Este capítulo presenta los resultados obtenidos al aplicar las
metodologías descritas en el capítulo anterior, así como del análisis
estadístico descriptivo de las variables estudiadas y las ventajas y
desventajas al aplicar dichos tests.\}

\section{Presentación de resultados 1}

Descripción\\

\section{Presentación de resultados 2}

Descripción\\

\section{Presentación de resultados n}

Descripción\\

\chapter{Conclusiones y Recomendaciones}

\textbf{1.-} Conclusion 1 del trabajo de investigación.\\

\textbf{Recomendación} Recomendación 1 del trabajo de investigación.\\

\textbf{2.-} Conclusion 2 del trabajo de investigación.\\

\textbf{Recomendación} Recomendación 2 del trabajo de investigación.\\

\textbf{n.-} Conclusion n del trabajo de investigación.\\

\textbf{Recomendación} Recomendación n del trabajo de investigación.\\

\appendix
\renewcommand{\thechapter}{\Alph{chapter}}

\chapter{Anexo A}

\chapter{Anexo B}

\chapter{Anexo C}

video-games \textless- read\_csv (``video\_games\_sales.csv'')

%%%%%%%%%%%%%%%%%%%%%%%%%%%%%%%%%%%%%%%%%%%%%%%%%%%%%
%                   REFERENCIAS                     %
%%%%%%%%%%%%%%%%%%%%%%%%%%%%%%%%%%%%%%%%%%%%%%%%%%%%%
% Mantenemos la estructura natbib de tu plantilla
\renewcommand{\bibname}{Lista de Referencias}
\bibliographystyle{apalike} 
\bibliography{bibliografia.bib} 

%%%%%%%%%%%%%%%%%%%%%%%%%%%%%%%%%%%%%%%%%%%%%%%%%%%%%
%                   APÉNDICES                       %
%%%%%%%%%%%%%%%%%%%%%%%%%%%%%%%%%%%%%%%%%%%%%%%%%%%%%

\end{document}
